\section{Introduction}
\label{sec:intro}

LLM agents are increasingly used for long-running tasks such as repository-level software engineering \citep{jimenez2024swebench,yang2024sweagent} and deep research over the open web \citep{zheng2025deepresearcher,wei2025browsecomp}.
Unlike single-turn generation, these tasks unfold over extended trajectories of model calls, tool executions, observations, failures, and revisions. As trajectories grow, a central challenge for the \emph{agent harness} is context management: session history accumulates continuously, while each model invocation operates over a bounded context window. Moreover, the effective context a model can reliably exploit is far smaller than its nominal window, as long-input retrieval and reasoning degrade with input length \citep{modarressi2025nolima,zeng2026loca}.

Current systems largely address this problem through \emph{context compression} or \emph{external memory}.
Compression is the dominant approach in practice: existing methods truncate stale spans, discard tool outputs, or replace earlier trajectory segments with summaries \citep{kang2025acon,ye2025agentfold}, and production agents such as Claude Code, Codex CLI, and Cursor reportedly employ similar compaction mechanisms as the context window approaches its limit.
External memory systems extract selected facts or episodes into a separate store and later retrieve them through semantic or structured interfaces \citep{packer2023memgpt,wu2025longmemeval, cao2026remember}. 
Both are fundamentally lossy in the same way: the agent sees history only through the summary or the memory store, so any detail they fail to keep is out of reach---even if the raw log still exists on disk.
Long-horizon tasks, however, may require exact historical evidence or nontrivial computation over past events, such as comparing tool outputs produced far apart in the trajectory. Neither the relevant information nor the required operation is known in advance, so no summary produced at observation time can be guaranteed to preserve what is later needed.


We present \emph{Scroll}, which keeps the agent's history outside the model context \citep{zhang2025rlm} and represents it as an executable \emph{Session Environment}. 
An append-only Event Log preserves the interaction trajectory with stable addresses and provenance, while a sandboxed, persistent Python kernel survives across model calls and maintains a typed namespace of resident variables. Any of this state can therefore be materialized as Python objects and reused across reasoning steps without being serialized into the prompt.

This turns context management into \emph{writing programs}---something current models are already highly proficient at. The model issues \texttt{exec} actions to search and expand the Event Log, access permitted resources, invoke tools \citep{ptc}, and compute over resident variables. Retrieved records, tool outputs, and intermediate computations remain in the kernel unless explicitly emitted through \texttt{print}, which the harness inserts as an observation into the next model context. Thus, \texttt{exec} determines how the environment is accessed and transformed, while \texttt{print} determines which projection enters the model's \emph{working view} over the Session Environment.

 As the working view approaches its budget, the harness evicts stale spans. Unlike compaction, eviction changes only the view, never the underlying record: evicted events remain verbatim in the Event Log under their stable addresses, where the model's programs can search for and materialize them on demand. 
Scroll additionally keeps an \emph{eviction index} in the view: a compact map of what has left it. Where search recovers only what the agent thinks to ask for, the index keeps the agent aware of history it can no longer see; each entry anchors the exact addresses of the evicted events, from which the originals are materialized on demand.

Our contributions are threefold:
\begin{itemize}[leftmargin=1.4em,itemsep=1pt,topsep=2pt]
    \item We formulate long-horizon context management as choosing, at each step, a working view over a persistent Session Environment. Existing approaches fix this choice before future needs are known; Scroll defers it to query time as a program the model writes.
    \item We implement an executable context substrate combining an append-only Event Log, durable storage, and a sandboxed persistent Python kernel. The model operates on the environment through \texttt{exec}; within it, only explicit \texttt{print} output crosses into the model-visible context.
    \item We introduce an algorithm that keeps the working view within budget without losing history: evicted spans remain intact in the Event Log, indexed by compact address-anchored entries the agent can navigate directly.
\end{itemize}